\begin{description}
\item[(a)] DES is Table 1 and SDSS Table 2, since faint stars will have
higher uncertainties. SDSS uses a smaller telescope (2.5\,m) than DES
(4\,m) hence, it is typically shallower and has larger errors. Just looking
at the error distribution should be enough to tell who has larger errors.

\item[(b)] The linear regression of each curve:

% FIGURE NEEDED: Graph 1 solution -- semi-log plot of err m_g vs m_g with
% DES (red crosses, dotted fit) and SDSS (blue plusses, solid fit) data and
% fitted lines (2024_da_sol.pdf page 2)

The reported parameters are:

\textbf{SDSS:}
\[ \log_{10}(err\ m_g) = A \cdot m_g + B \]
$A = 0.2780$ (acceptable from 0.2502 to 0.3058 which is within 10\% range
of the fit)\\
$B = -7.3108$ (acceptable from $-8.0419$ to $-6.5797$ which is within 10\%
range of the fit)

\textbf{DES:}
\[ \log_{10}(err\ m_g) = A \cdot m_g + B \]
$A = 0.4063$ (acceptable from 0.3657 to 0.4470 which is within 10\% range
of the fit)\\
$B = -10.7597$ (acceptable from $-10.2217$ to $-11.2977$ which is within
10\% range of the fit)

\item[(c)] Using that $S/N \sim 1/err\ m_g$ and above fits we can arrive at
the following:\\
\textbf{SDSS:} $\log_{10}(err\ m_g) = 0.2780 \cdot 21.5 - 7.3108
\rightarrow err\ m_g = 0.0464$\\
\textbf{DES:} $\log_{10}(err\ m_g) = 0.4063 \cdot 21.5 - 10.7597
\rightarrow err\ m_g = 0.0095$\\
Accepted answers follow below.

\begin{center}
\begin{tabular}{|l|c|l|}
\hline
 & err & $S/N$ \\
\hline
SDSS err at 21.5 mag & 0.05 & \textbf{22} (acceptable from 19.8 to 24.2
$\rightarrow$ within 10\% range of the result) \\
\hline
DES err at 21.5 mag & 0.01 & \textbf{106} (acceptable from 95.4 to 116.6
$\rightarrow$ within 10\% range of the result) \\
\hline
\end{tabular}
\end{center}

\item[(d)] One tip for this question is that the student should realize
that the SDSS RA coordinate is sorted, so it can be used as a reference for
the scanning of DES coordinates to perform the match. These are the stars
that can be matched between catalogs:

\begin{center}
\begin{tabular}{ccccc|ccccc|c}
$ID_1$ & RA & Dec & $m_g$ & err $m_g$ &
$ID_2$ & RA & Dec & $m_g$ & err $m_g$ & Sep \\
 & (deg) & (deg) & (mag) & (mag) &
 & (deg) & (deg) & (mag) & (mag) & (arcsec) \\
\hline
3 & 0.064911 & 0.026395 & 21.3931 & 0.0084 & 8 & 0.064910 & 0.026419 & 21.6428 & 0.0488 & 0.08475 \\
9 & 0.057422 & 0.076006 & 21.8897 & 0.0127 & 6 & 0.057414 & 0.075999 & 21.8884 & 0.0578 & 0.03724 \\
4 & 0.098343 & 0.054871 & 21.3934 & 0.0088 & 13 & 0.098343 & 0.054854 & 21.6542 & 0.0499 & 0.06186 \\
6 & 0.006188 & 0.066928 & 21.5490 & 0.0088 & 1 & 0.006167 & 0.066874 & 21.9020 & 0.0576 & 0.2076 \\
12 & 0.071089 & 0.091053 & 21.4390 & 0.0098 & 9 & 0.071102 & 0.091058 & 21.9259 & 0.0751 & 0.05009 \\
2 & 0.064741 & 0.021568 & 21.1111 & 0.0067 & 7 & 0.064745 & 0.021583 & 21.3634 & 0.0422 & 0.05655 \\
8 & 0.076715 & 0.079496 & 21.4808 & 0.0089 & 11 & 0.076709 & 0.079747 & 21.5303 & 0.0476 & 0.08276 \\
% sic: the solution table gives Dec 0.079747 for ID2 = 11, whereas Table 2
% of the problem lists 0.079474
\end{tabular}
\end{center}

\item[(e)] The figure of g magnitude of DES vs SDSS should look like this
with the best linear fit:

% FIGURE NEEDED: Graph 2 solution -- SDSS magnitude vs DES magnitude with
% 2-sigma error bars and the one-to-one line "suitable for photometric
% calibration (x = y)" (2024_da_sol.pdf page 3), plus a second version with
% a non-zero-offset line (2024_da_sol.pdf page 4)

The student should identify the stars that are closest to the one-to-one
line ($x = y$) between surveys. This could be with a zero offset, or a
non-zero offset. The \textbf{two} stars that would be suitable for
photometric calibration with zero offset are the ones with ID's \textbf{8}
and \textbf{9} of Table 1 (DES). The \textbf{four} stars that would be
suitable for photometric calibration with non-zero offset are the ones with
ID's \textbf{2, 3, 4 and 6} of Table 1 (DES).
\end{description}
